\IfFileExists{stacks-project.cls}{%
\documentclass{stacks-project}
}{%
\documentclass{amsart}
}
\usepackage{amssymb}

% For dealing with references we use the comment environment
\usepackage{verbatim}
\usepackage{fancyvrb}
\usepackage{pmboxdraw}
\newenvironment{reference}{\comment}{\endcomment}
%\newenvironment{reference}{}{}
\newenvironment{slogan}{\comment}{\endcomment}
\newenvironment{history}{\comment}{\endcomment}

% For source figures retained as assets
\usepackage{graphicx}

% For commutative diagrams we use Xy-pic
\usepackage[all]{xy}

% We use 2cell for 2-commutative diagrams.
\xyoption{2cell}
\UseAllTwocells

% We use multicol for the list of chapters between chapters
\usepackage{multicol}

% This is generally recommended for better output
\usepackage{lmodern}
\usepackage[T1]{fontenc}
\usepackage{textalpha}

% For cross-file-references
% We prefer the snapshot tmp/xref/#2.aux,
% which latexmkrc replaces atomically, so
% that documents compiled simultaneously
% never read a half-written .aux file.
% The last argument keeps links to #2.pdf.
\usepackage{xr-hyper}
\newcommand{\maybeexternaldocument}[2]{%
  \IfFileExists{#2.aux}{%
    \externaldocument[#1]{#2}%
  }{%
    \PackageWarning{stack}{Missing #2.aux; external references with prefix #1 unavailable}%
  }%
}
\let\externaldocumentorig\externaldocument
\renewcommand{\externaldocument}[2][]{%
  \IfFileExists{tmp/xref/#2.aux}{%
    \externaldocumentorig[#1]{tmp/xref/#2}[#2.pdf]%
  }{%
  \IfFileExists{#2.aux}{%
    \externaldocumentorig[#1]{#2}%
  }{%
    \PackageWarning{stack}{Missing #2.aux; external references with prefix #1 unavailable}%
  }%
  }%
}

% Package for hypertext links:
\usepackage{hyperref}

% For any local file, say "hello.tex" you want to link to please
% use \externaldocument[hello-]{hello}
\externaldocument[introduction-]{introduction}
\externaldocument[conventions-]{conventions}
\externaldocument[sets-]{sets}
\externaldocument[categories-]{categories}
\externaldocument[topology-]{topology}
\externaldocument[sheaves-]{sheaves}
\externaldocument[sites-]{sites}
\externaldocument[stacks-]{stacks}
\externaldocument[fields-]{fields}
\externaldocument[algebra-]{algebra}
\externaldocument[brauer-]{brauer}
\externaldocument[homology-]{homology}
\externaldocument[derived-]{derived}
\externaldocument[simplicial-]{simplicial}
\externaldocument[more-algebra-]{more-algebra}
\externaldocument[smoothing-]{smoothing}
\externaldocument[modules-]{modules}
\externaldocument[sites-modules-]{sites-modules}
\externaldocument[injectives-]{injectives}
\externaldocument[cohomology-]{cohomology}
\externaldocument[sites-cohomology-]{sites-cohomology}
\externaldocument[dga-]{dga}
\externaldocument[dpa-]{dpa}
\externaldocument[sdga-]{sdga}
\externaldocument[hypercovering-]{hypercovering}
\externaldocument[schemes-]{schemes}
\externaldocument[constructions-]{constructions}
\externaldocument[properties-]{properties}
\externaldocument[morphisms-]{morphisms}
\externaldocument[coherent-]{coherent}
\externaldocument[divisors-]{divisors}
\externaldocument[limits-]{limits}
\externaldocument[varieties-]{varieties}
\externaldocument[topologies-]{topologies}
\externaldocument[descent-]{descent}
\externaldocument[perfect-]{perfect}
\externaldocument[more-morphisms-]{more-morphisms}
\externaldocument[flat-]{flat}
\externaldocument[groupoids-]{groupoids}
\externaldocument[more-groupoids-]{more-groupoids}
\externaldocument[etale-]{etale}
\externaldocument[chow-]{chow}
\externaldocument[intersection-]{intersection}
\externaldocument[pic-]{pic}
\externaldocument[weil-]{weil}
\externaldocument[adequate-]{adequate}
\externaldocument[dualizing-]{dualizing}
\externaldocument[duality-]{duality}
\externaldocument[discriminant-]{discriminant}
\externaldocument[derham-]{derham}
\externaldocument[local-cohomology-]{local-cohomology}
\externaldocument[algebraization-]{algebraization}
\externaldocument[curves-]{curves}
\externaldocument[resolve-]{resolve}
\externaldocument[models-]{models}
\externaldocument[functors-]{functors}
\externaldocument[equiv-]{equiv}
\externaldocument[pione-]{pione}
\externaldocument[etale-cohomology-]{etale-cohomology}
\externaldocument[proetale-]{proetale}
\externaldocument[relative-cycles-]{relative-cycles}
\externaldocument[more-etale-]{more-etale}
\externaldocument[trace-]{trace}
\externaldocument[crystalline-]{crystalline}
\externaldocument[spaces-]{spaces}
\externaldocument[spaces-properties-]{spaces-properties}
\externaldocument[spaces-morphisms-]{spaces-morphisms}
\externaldocument[decent-spaces-]{decent-spaces}
\externaldocument[spaces-cohomology-]{spaces-cohomology}
\externaldocument[spaces-limits-]{spaces-limits}
\externaldocument[spaces-divisors-]{spaces-divisors}
\externaldocument[spaces-over-fields-]{spaces-over-fields}
\externaldocument[spaces-topologies-]{spaces-topologies}
\externaldocument[spaces-descent-]{spaces-descent}
\externaldocument[spaces-perfect-]{spaces-perfect}
\externaldocument[spaces-more-morphisms-]{spaces-more-morphisms}
\externaldocument[spaces-flat-]{spaces-flat}
\externaldocument[spaces-groupoids-]{spaces-groupoids}
\externaldocument[spaces-more-groupoids-]{spaces-more-groupoids}
\externaldocument[bootstrap-]{bootstrap}
\externaldocument[spaces-pushouts-]{spaces-pushouts}
\externaldocument[spaces-chow-]{spaces-chow}
\externaldocument[groupoids-quotients-]{groupoids-quotients}
\externaldocument[spaces-more-cohomology-]{spaces-more-cohomology}
\externaldocument[spaces-simplicial-]{spaces-simplicial}
\externaldocument[spaces-duality-]{spaces-duality}
\externaldocument[formal-spaces-]{formal-spaces}
\externaldocument[restricted-]{restricted}
\externaldocument[spaces-resolve-]{spaces-resolve}
\externaldocument[formal-defos-]{formal-defos}
\externaldocument[defos-]{defos}
\externaldocument[cotangent-]{cotangent}
\externaldocument[examples-defos-]{examples-defos}
\externaldocument[algebraic-]{algebraic}
\externaldocument[examples-stacks-]{examples-stacks}
\externaldocument[stacks-sheaves-]{stacks-sheaves}
\externaldocument[criteria-]{criteria}
\externaldocument[artin-]{artin}
\externaldocument[quot-]{quot}
\externaldocument[stacks-properties-]{stacks-properties}
\externaldocument[stacks-morphisms-]{stacks-morphisms}
\externaldocument[stacks-limits-]{stacks-limits}
\externaldocument[stacks-cohomology-]{stacks-cohomology}
\externaldocument[stacks-perfect-]{stacks-perfect}
\externaldocument[stacks-introduction-]{stacks-introduction}
\externaldocument[stacks-more-morphisms-]{stacks-more-morphisms}
\externaldocument[stacks-geometry-]{stacks-geometry}
\externaldocument[moduli-]{moduli}
\externaldocument[moduli-curves-]{moduli-curves}
\externaldocument[examples-]{examples}
\externaldocument[exercises-]{exercises}
\externaldocument[guide-]{guide}
\externaldocument[desirables-]{desirables}
\externaldocument[coding-]{coding}
\externaldocument[obsolete-]{obsolete}
\externaldocument[fdl-]{fdl}
\externaldocument[index-]{index}

% Heap Project documents
\externaldocument[k3-]{k3}
\externaldocument[torelli-]{torelli}
\externaldocument[linear-algebra-]{linear-algebra}
\externaldocument[groups-]{groups}
\externaldocument[complex-analysis-]{complex-analysis}
\externaldocument[functional-analysis-]{functional-analysis}
\externaldocument[categories-]{categories}
\externaldocument[commutative-algebra-]{commutative-algebra}
\externaldocument[ringed-spaces-]{ringed-spaces}
\externaldocument[homological-algebra-]{homological-algebra}
\externaldocument[lie-algebras-]{lie-algebras}
\externaldocument[representation-theory-]{representation-theory}
\externaldocument[smooth-manifolds-]{smooth-manifolds}
\externaldocument[differential-topology-]{differential-topology}
\externaldocument[algebraic-topology-]{algebraic-topology}
\externaldocument[differential-geometry-]{differential-geometry}
\externaldocument[algebraic-geometry-]{algebraic-geometry}
\externaldocument[symplectic-geometry-]{symplectic-geometry}
\externaldocument[complex-geometry-]{complex-geometry}
\externaldocument[hodge-theory-]{hodge-theory}
\externaldocument[%
variations-of-hodge-structures-]{%
variations-of-hodge-structures}
\externaldocument[%
homological-mirror-symmetry-]{%
homological-mirror-symmetry}
\maybeexternaldocument{deformations-}{deformations}
\externaldocument[vertex-algebras-]{vertex-algebras}
\externaldocument[springer-]{springer}
\externaldocument[infty-categories-]{infty-categories}
\externaldocument[seminars-differential-geometry-]{%
seminars-differential-geometry}
\externaldocument[seminars-complex-geometry-]{%
seminars-complex-geometry}
\externaldocument[seminars-algebraic-geometry-]{%
seminars-algebraic-geometry}
\externaldocument[seminars-analysis-pde-]{%
seminars-analysis-pde}
\externaldocument[surfaces-]{surfaces}
\externaldocument[birational-maps-commutative-algebra-]{birational-maps-commutative-algebra}
\externaldocument[cremona-transformations-]{cremona-transformations}
\externaldocument[derived-categories-of-sheaves-]{derived-categories-of-sheaves}
\externaldocument[geometric-invariant-theory-]{geometric-invariant-theory}
\externaldocument[geometric-stability-conditions-and-group-actions-]{geometric-stability-conditions-and-group-actions}
\externaldocument[moduli-spaces-of-sheaves-]{moduli-spaces-of-sheaves}
\externaldocument[bridgeland-stability-general-theory-]{bridgeland-stability-general-theory}
\externaldocument[hilbert-schemes-]{hilbert-schemes}
\externaldocument[noncommutative-abelian-surfaces-and-kummer-type-hyperkahler-varieties-]{noncommutative-abelian-surfaces-and-kummer-type-hyperkahler-varieties}
\externaldocument[mumford-tate-groups-in-hodge-theory-]{mumford-tate-groups-in-hodge-theory}
\externaldocument[hodge-birational-atoms-2026-]{hodge-birational-atoms-2026}
\externaldocument[xxii-egd-2026-]{%
xxii-egd-2026}
\externaldocument[pde-]{pde}
\externaldocument[probability-]{probability}

\externaldocument[linguistics-]{linguistics}

\externaldocument[%
16th-alga-meeting-2026-]{%
16th-alga-meeting-2026}

\externaldocument[%
iv-jornadas-lefschetz-2026-]{%
iv-jornadas-lefschetz-2026}

\externaldocument[reading-the-masters-]{%
reading-the-masters}

\externaldocument[real-analysis-]{%
real-analysis}

\externaldocument[optimal-transport-]{%
optimal-transport}

\externaldocument[history-]{history}

\externaldocument[seiberg-witten-]{seiberg-witten}

\externaldocument[history-philosophy-biology-]{history-philosophy-biology}

\externaldocument[classical-mechanics-]{classical-mechanics}
\externaldocument[electromagnetism-]{electromagnetism}
\externaldocument[waves-and-optics-]{waves-and-optics}
\externaldocument[wif-]{waves-in-fluids-iv}
\externaldocument[thermal-physics-]{thermal-physics}
\externaldocument[quantum-mechanics-]{quantum-mechanics}
\externaldocument[special-relativity-]{special-relativity}
\externaldocument[physics-]{physics}

\externaldocument[chemistry-]{chemistry}

\externaldocument[molecular-cell-biology-]{molecular-cell-biology}

\externaldocument[microbiology-]{microbiology}

\externaldocument[protistology-]{protistology}

\externaldocument[mycology-]{mycology}

\externaldocument[botany-]{botany}

\externaldocument[zoology-]{zoology}

\externaldocument[animal-histology-]{animal-histology}

\externaldocument[animal-physiology-]{animal-physiology}

\externaldocument[systematics-biogeography-]{systematics-biogeography}

\externaldocument[ecology-natural-resources-]{ecology-natural-resources}

\externaldocument[earth-sciences-]{earth-sciences}

\externaldocument[biostatistics-]{biostatistics}

\externaldocument[mathematical-systems-biology-]{mathematical-systems-biology}

\externaldocument[cimpa-el-salvador-2024-]{cimpa-el-salvador-2024}

\externaldocument[anthropology-]{anthropology}

\externaldocument[yang-mills-theory-]{yang-mills-theory}

\externaldocument[seminars-algebra-]{seminars-algebra}

\externaldocument[seminars-physics-]{seminars-physics}

\externaldocument[seminars-computing-]{seminars-computing}

\externaldocument[oktobermat-2026-]{oktobermat-2026}

\externaldocument[saet-2026-]{saet-2026}

\externaldocument[seminars-symplectic-geometry-]{seminars-symplectic-geometry}

\externaldocument[seminars-dynamics-]{seminars-dynamics}

\externaldocument[oktobermat-2024-]{oktobermat-2024}

\externaldocument[quantum-groups-]{quantum-groups}

\externaldocument[legal-2027-]{legal-2027}

\externaldocument[derived-categories-]{derived-categories}

% Theorem environments.
%
\theoremstyle{plain}
\newtheorem{theorem}[subsection]{Theorem}
\newtheorem{proposition}[subsection]{Proposition}
\newtheorem{lemma}[subsection]{Lemma}
\newtheorem{corollary}[subsection]{Corollary}

\theoremstyle{definition}
\newtheorem{definition}[subsection]{Definition}
\newtheorem{example}[subsection]{Example}
\newtheorem{exercise}[subsection]{Exercise}
\newtheorem{situation}[subsection]{Situation}

\theoremstyle{remark}
\newtheorem{remark}[subsection]{Remark}
\newtheorem{remarks}[subsection]{Remarks}

\numberwithin{equation}{subsection}

% Macros
%
\def\lim{\mathop{\mathrm{lim}}\nolimits}
\def\colim{\mathop{\mathrm{colim}}\nolimits}
\def\Spec{\mathop{\mathrm{Spec}}}
\def\Hom{\mathop{\mathrm{Hom}}\nolimits}
\def\Ext{\mathop{\mathrm{Ext}}\nolimits}
\def\SheafHom{\mathop{\mathcal{H}\!\mathit{om}}\nolimits}
\def\SheafExt{\mathop{\mathcal{E}\!\mathit{xt}}\nolimits}
\def\Sch{\mathit{Sch}}
\def\Mor{\mathop{\mathrm{Mor}}\nolimits}
\def\Ob{\mathop{\mathrm{Ob}}\nolimits}
\def\Sh{\mathop{\mathit{Sh}}\nolimits}
\def\NL{\mathop{N\!L}\nolimits}
\def\CH{\mathop{\mathrm{CH}}\nolimits}
\def\proetale{{pro\text{-}\acute{e}tale}}
\def\etale{{\acute{e}tale}}
\def\QCoh{\mathit{QCoh}}
\def\Ker{\mathop{\mathrm{Ker}}}
\def\Im{\mathop{\mathrm{Im}}}
\def\Coker{\mathop{\mathrm{Coker}}}
\def\Coim{\mathop{\mathrm{Coim}}}

% Boxtimes
%
\DeclareMathSymbol{\boxtimes}{\mathbin}{AMSa}{"02}

%
% Macros for moduli stacks/spaces
%
\def\QCohstack{\mathcal{QC}\!\mathit{oh}}
\def\Cohstack{\mathcal{C}\!\mathit{oh}}
\def\Spacesstack{\mathcal{S}\!\mathit{paces}}
\def\Quotfunctor{\mathrm{Quot}}
\def\Hilbfunctor{\mathrm{Hilb}}
\def\Curvesstack{\mathcal{C}\!\mathit{urves}}
\def\Polarizedstack{\mathcal{P}\!\mathit{olarized}}
\def\Complexesstack{\mathcal{C}\!\mathit{omplexes}}
% \Pic is the operator that assigns to X its picard group, usage \Pic(X)
% \Picardstack_{X/B} denotes the Picard stack of X over B
% \Picardfunctor_{X/B} denotes the Picard functor of X over B
\def\Pic{\mathop{\mathrm{Pic}}\nolimits}
\def\Picardstack{\mathcal{P}\!\mathit{ic}}
\def\Picardfunctor{\mathrm{Pic}}
\def\Deformationcategory{\mathcal{D}\!\mathit{ef}}


\begin{document}

\title{Quantum groups}
\maketitle

\phantomsection
\label{section-phantom}

\tableofcontents

\noindent
These are Daniel's own notes of a
mini-course on quantum groups by Artem,
in the summer of 2024. There were four
lectures; the first one took place on
July 11, 2024, from 15:00 to 16:30.
The surname of the lecturer and the event
where the course was given are not
recorded in the notes.

\medskip\noindent
{\bf Editorial note.}
In the original notes, places where
Daniel was unsure of what was written on
the board were marked in magenta; here
these marks are removed and the question
marks or gaps they marked are kept as
plain text. Pictures and tikz-cd diagrams
were replaced by descriptions or by
xymatrix diagrams. Obvious misspellings
(for instance of Verma, Yangian, Yang,
Baxter, Reidemeister, Belavin, Whitehead)
were corrected. Other repairs are listed
at the end of each lecture.

\section{First lecture: Hopf algebras}
\label{section-first-lecture}

\noindent
Daniel's summary of the lecture, as
written at the top of the notes:
``On July 11 from 15:00 to 16:30 we had
a first lecture of a mini-course on
quantum groups by Artem. He defined Hopf
algebras, gave some examples: functions
on algebraic group, dual, universal
enveloping algebra for a Lie algebra,
particular examples with $g^n=1$,
$x^n=0$, $gx = z xg$ with
$z = \exp(2\pi i/n)$. Reminder on
representations of $\mathfrak{sl}_{2}$ and
Verma modules.''

\medskip\noindent
``First lecture was finished with a
quantum cliffhanger: the universal
enveloping algebra
$\operatorname{U}(\mathfrak{g}) =
\operatorname{U}_{\hbar}(\mathfrak{g})$
seen as a (deformation) quantization of
$\mathfrak{g}^*$, the latter considered
as a Poisson variety (after Sophus Lie
and KKS).''

\medskip\noindent
{\bf Definition of Hopf algebras.}
$G$ is an affine algebraic variety over
$\mathbb{C}$. (A group object in the
category of algebraic varieties.)
We shall work mostly with
$\operatorname{GL}_{n}$ and
$\operatorname{SL}_{n}$. There are also
finite groups, $\operatorname{SO}_{n}$,
$\operatorname{Sp}_{n}$. So $G$ is a
linear group.

Everything we say of $G$ will be
determined by its ring of functions
$\mathcal{O}(G)=H$, which is an algebra.

We have the {\it multiplication map} and
the {\it unit map}
$$
m:H\otimes H\to H
\qquad\qquad
u:\mathbb{C}\to H
$$
for which there are some axioms:
\begin{enumerate}
\item[{\bf A1}] Associativity. The
following diagram commutes:
$$
\xymatrix{
H\otimes H\otimes H
\ar[r]^-{m\otimes\operatorname{id}}
\ar[d]_{\operatorname{id}\otimes m}
& H\otimes H \ar[d]^{m}\\
H\otimes H \ar[r]_-{m} & H
}
$$
\item[{\bf A2}] Unitality.
$$
\xymatrix{
H\cong H\otimes\mathbb{C}
\ar[r]^-{\operatorname{id}\otimes u}
\ar@/_1.5pc/[rr]_{\operatorname{id}}
& H\otimes H \ar[r]^-{m} & H
}
$$
\end{enumerate}

\begin{definition}
\label{definition-l1-algebra}
An {\it algebra} is a structure with
multiplication and unit map subject to
axioms A1 and A2.
\end{definition}

Now comes the group part. So what does it
mean to be a group? That we have a map
$G\times G\to G$ and a unit
$\mathbb{C}\to G$. On the level of
functions these correspond to the
{\it coproduct}
$$
\Delta:H\to H\otimes H
$$
and the {\it counit}
$$
\eta:H\to \mathbb{C}
$$

And we put some new axioms:
\begin{enumerate}
\item[{\bf A3}] Coassociativity.
$$
\xymatrix{
H \ar[r]^-{\Delta} \ar[d]_{\Delta}
& H\otimes H
\ar[d]^{\Delta\otimes\operatorname{id}}\\
H\otimes H
\ar[r]_-{\operatorname{id}\otimes\Delta}
& H\otimes H\otimes H
}
$$
\item[{\bf A4}] Counitality.
$$
\xymatrix{
H \ar[r]^-{\Delta}
\ar@/_1.5pc/[rr]_{\operatorname{id}}
& H\otimes H
\ar[r]^-{\eta\otimes\operatorname{id}} & H
}
$$
\end{enumerate}

\begin{definition}
\label{definition-l1-coalgebra}
A {\it coalgebra} is a structure with a
coproduct and a counit subject to
axioms A3 and A4.
\end{definition}

And then we ask that the two structures
are compatible with the following axiom:
\begin{enumerate}
\item[{\bf A5}]
$$
\xymatrix{
H\otimes H \ar[rr]^-{m}
\ar[d]_{\Delta\otimes\Delta}
&& H \ar[d]^{\Delta}\\
H\otimes H\otimes H\otimes H
\ar[dr]_{\operatorname{id}\otimes P
\otimes\operatorname{id}}
&& H\otimes H\\
& H\otimes H\otimes H\otimes H
\ar[ur]_{m\otimes m}
}
$$
\end{enumerate}

\begin{remark}
\label{remark-l1-homomorphisms}
So we are asking that coproduct and
product are homomorphisms of algebras.
And there are also some properties about
the units.
\end{remark}

\begin{remark}
\label{remark-l1-other-axioms}
There are other axioms we did not write
down. One of them says that some map is a
map of $k$-algebras. Anyway,
\end{remark}

\begin{definition}
\label{definition-l1-bialgebra}
A {\it bialgebra} is a structure with a
multiplication, unit, coproduct and
counit satisfying axioms A1 to A8.
\end{definition}

So in a group we have the map
$G\to G$ mapping $g\mapsto g^{-1}$, which
corresponds to the {\it antipode map}
$$
S:H\to H
$$

Let us write this other axiom:
\begin{enumerate}
\item[{\bf A9}]
$$
\xymatrix{
& H\otimes H
\ar[rr]^-{S\otimes\operatorname{id}}
&& H\otimes H \ar[dr]^{m}\\
H \ar[ur]^{\Delta} \ar[dr]_{\Delta}
\ar[rr]^-{\eta}
&& \mathbb{C} \ar[rr]^-{u} && H\\
& H\otimes H
\ar[rr]_-{\operatorname{id}\otimes S}
&& H\otimes H \ar[ur]_{m}
}
$$
\end{enumerate}

\begin{definition}
\label{definition-l1-hopf-algebra}
A {\it Hopf algebra} is a structure with
multiplication, unit, product, coproduct
and antipode map satisfying axioms
A1-A9.
\end{definition}

\begin{remark}
\label{remark-l1-commutative}
For $\mathcal{O}(G)$, $m$ is commutative:
$$
\xymatrix{
H\otimes H \ar[rr]^-{m}
\ar[dr]_{P} && H\\
& H\otimes H \ar[ur]_{m}
}
$$
This is a special feature of affine
algebraic groups.
\end{remark}

\begin{definition}
\label{definition-l1-cocommutative}
The coproduct $\Delta$ is
{\it cocommutative} if
$$
\xymatrix{
H \ar[rr]^-{\Delta}
\ar[dr]_{\Delta} && H\otimes H\\
& H\otimes H \ar[ur]_{P}
}
$$
commutes.
\end{definition}

\begin{example}[1, Functions on algebraic
groups]
\label{example-l1-functions}
Let $G$ be any affine algebraic group like
$\operatorname{GL}_{n}$, the invertible
matrices. What is the ring of functions?
Let us present it in a flappy way:
$$
\mathcal{O}(\operatorname{GL}_{n})=
\mathbb{C}\left\langle{}
t_{ij}\,|\,t_{ij}=t_{ji}
\right\rangle{}
$$
where $t_{ij}$ are the coordinates of a
matrix, with the following relations:
$$
T=(t_{ij})\in
\operatorname{End}(\mathbb{C}^{n})
\otimes H
$$
so that
\begin{itemize}
\item $T_{1}T_{2}=T_{2}T_{1}$
which means there is a commutativity
condition: $t_{ij}t_{ik}=t_{ik}t_{ij}$.

\item (Invertibility of $\det$.)
$\mathcal{O}(\operatorname{GL}_{n})=
\mathbb{C}[t_{ij},\det{}^{-1}]$

So that
$\mathcal{O}(\operatorname{SL}_{n})=
\mathbb{C}[t_{ij}]/(\det =1)$
\end{itemize}
The coproduct here is
$$
\Delta T=T\otimes T
$$
so
$$
\Delta t_{ij}=\sum_{k=1}^{n}
t_{ik}\otimes t_{kj}
$$
which is what matrix multiplication does.

And the counit sends $T$ to the identity
matrix, that is,
$$
\eta(t_{ij})=\delta_{ij}
$$

And the antipode is
$$
S(T)=T^{-1}
$$
\end{example}

\begin{example}[2, Dual]
\label{example-l1-dual}
If $H$ is a finite-dimensional Hopf
algebra, its dual $H^{*}$ obtained by
reversing all arrows in the construction
of $H$ is also a Hopf algebra.

If $G$ is a finite group and
$H=\mathcal{O}(G)$ then
$H^{*} =\mathbb{C}[G]$ is a
{\it group algebra}. Multiplication is the
usual multiplication.

Another example is if $G$ is a connected
affine algebraic group. Then
$\mathcal{O}(G)=H$ is infinite
dimensional, but there is kind of the
dual, which is called the
{\it universal enveloping algebra}
$U\mathfrak{g}$.

Recall that a Lie algebra is
$\mathfrak{g}=T_{e}G$. For example, if
$G=\operatorname{GL}_{n}$ then
$\mathfrak{g}=
\operatorname{End}(\mathbb{C}^{n})=
\mathfrak{gl}_{n}$. Abstractly, a Lie
algebra $\mathfrak{g}$ is a vector space
with
$[\,,]:\mathfrak{g}\otimes\mathfrak{g}
\to\mathfrak{g}$ such that
$[x,y]=-[y,x]$ and
$[x,[y,z]]+[y,[x,z]]+[z,[x,y]]=0$.
For example, $\mathfrak{gl}_{n}$ with
bracket $[A,B]=AB-BA$.

Then for the universal enveloping algebra
take the tensor algebra
$$
T(V)=\bigoplus_{k} V^{\otimes k}
$$
then
$$
U\mathfrak{g}=T(\mathfrak{g})/(xy-yx-[x,y])
$$
which turns out to be a Hopf algebra for
$x\in \mathfrak{g}$ with
\begin{itemize}
\item Coproduct: $\Delta(x)=x\otimes\ ?$
\item Counit: $\eta(x)=\ ?$
\item Antipode: $S(x)=-x$.
\end{itemize}
There is a kind of duality or
non-degenerate pairing
$$
U\mathfrak{g}\otimes \mathcal{O}(G)
\to \mathbb{C}
$$
\end{example}

\begin{example}[3, exercise]
\label{example-l1-exercise}
Let $H=\mathbb{C}\left\langle{}
g,x\right\rangle{} / (g^{n}=1,\ x^{n}=0,\
gx=\zeta xg)$ for $\zeta\in \mathbb{C}$
such that $\zeta^{n}=1$.

There is a coproduct given by
$$
\begin{aligned}
\Delta(g)&=g\otimes g\\
\Delta(x)&=x\otimes\ ?+g\otimes x
\end{aligned}
$$
\begin{itemize}
\item Show that this is a Hopf algebra.
In particular, find $\eta$ and $\zeta$.
\item Show that $\dim H=n^{2}$.
\item {\bf Hint:} use representation of
$H$.
\end{itemize}
\end{example}

Let $H$ be an algebra. Then
$\operatorname{Rep}H$ is the ``category''
of algebraic representations. Let $H$ be
in addition a bialgebra, then for any
$V,W\in \operatorname{Rep}H$,
$$
V\otimes W\in \operatorname{Rep}H
$$
and this works as follows: for
$v\in V$, $w\in W$ and $h \in H$,
$$
h(v\otimes w)=\Delta h\cdot
(v\otimes w)\in H\otimes H
$$
If additionally $H$ is a Hopf algebra,
i.e. there is an antipode $S$, we can
take duals:
$$
V\in \operatorname{Rep}H\implies
V^{*}\in \operatorname{Rep}H
$$
so, for a functional $\varphi \in V^{*}$
and $x\in V$, $h\in H$,
$(h\varphi)(x)=\varphi (S(h)x)$.

\begin{example}
\label{example-l1-group-algebra}
$H=\mathbb{C}[G]$,
$(g\cdot \varphi)(x)=\varphi (g^{-1}x)$.
\end{example}

\begin{remark}
\label{remark-l1-permutation}
As vector spaces:
$$
\begin{aligned}
P: V\otimes W &\longrightarrow
W\otimes V \\
v\otimes w &\longmapsto w\otimes v
\end{aligned}
$$
but it is not an isomorphism of
$H$-modules in general. If $H$ is
cocommutative, like $U\mathfrak{g}$, for
example, it is.
\end{remark}

\medskip\noindent
{\bf Review of representation theory of
$\mathfrak{sl}_{2}$.}
Let $\mathfrak{g}=\mathfrak{sl}_{2}\subset
\mathfrak{gl}_{2}$, the traceless
matrices. They are generated by
$$
e=\begin{pmatrix} 0&1\\0&0 \end{pmatrix},
\qquad
h=\begin{pmatrix} 1&0\\0&-1 \end{pmatrix},
\qquad
f=\begin{pmatrix} 0&0\\1&0 \end{pmatrix}
$$
with the relations
$$
[h,e]=2e,\qquad [h,f]=-2f,
\qquad [e,f]=h.
$$
There are some representations:
\begin{itemize}
\item (Trivial.) $V(0)=
\operatorname{span}(v_{0})$ with
$e(v_{0})=hv_{0}=fv_{0}=0$.
\item (Vector or fundamental
representation.) $V(1)=
\operatorname{span}(v_{0},v_{1})$, with an
action by matrices.
\end{itemize}

\begin{theorem}
\label{theorem-l1-sl2-representations}
\begin{enumerate}
\item Every finite dimensional
representation of $\mathfrak{sl}_{2}$ is
completely reducible,
$$
V=\bigoplus_{i}V_{i}
$$
where $V_{i}$ is irreducible.
\item Every irreducible is isomorphic to
$V(k)=\operatorname{span}
(v_{0},\ldots,v_{k})$ with
$$
\begin{aligned}
hv_{n}&=(\lambda-2n)v_{n}\\
ev_{n}&=(\lambda-?)\\
fv_{n}&=\ ?
\end{aligned}
$$
As an exercise, it can be shown that
$V(k)\cong \operatorname{Sym} V\ ?$.
\end{enumerate}
\end{theorem}

Now let us denote $\mathfrak{b}=
\operatorname{span}(h,e)\subset
\mathfrak{sl}_{2}$, so that
$U\mathfrak{b}\subset U\mathfrak{g}$. So
there is a representation of
$U\mathfrak{b}$ acting on
$\mathbb{C}_{\lambda}=
\operatorname{span}(v_{0}^{\lambda})$
and $hv_{0}^{\lambda} =\lambda
v_{0}^{\lambda}$ for
$\lambda\in \mathbb{C}$ and
$ev_{0}^{\lambda}=0$.

\begin{definition}
\label{definition-l1-verma-module}
The {\it Verma module} of weight
$\lambda$ is
$$
U\mathfrak{g}\otimes_{U\mathfrak{b}}
\mathbb{C}_{\lambda}.
$$
Concretely, as a vector space,
$$
M_{\lambda}=\operatorname{span}
(v_{0}^{\lambda},\
v_{1}^{\lambda} =fv_{0}^{\lambda},\
v_{2}^{\lambda}
=\tfrac{f^{2}}{2!} v_{0}^{\lambda},
\ldots)
$$
with the action
$$
\begin{aligned}
h\cdot v_{n}^{\lambda}
&=(\lambda-2n)v_{n}^{\lambda} \\
e\cdot v_{n}^{\lambda}
&=(\lambda-n+1)v_{n-1}^{\lambda}\\
f \cdot v_{n}^{\lambda}
&=(n+1)v^{\lambda}_{n+1}
\end{aligned}
$$
\end{definition}

\begin{proposition}
\label{proposition-l1-quotient}
If $\lambda\in \mathbb{Z}_{\geq 1}$, then
$V(\lambda)$ is a quotient of
$M_{\lambda}$.
\end{proposition}

Some terminology:
\begin{itemize}
\item $\lambda$ is the {\it highest
weight},
\item $h\cdot v_{n}^{\lambda}
=(\lambda-2n)v_{n}^{\lambda}$ is the
{\it weight} of $v_{n}^{\lambda}$
\item $v_{0}^{\lambda}$ is the
{\it highest weight vector}.
\end{itemize}

\medskip\noindent
{\bf Intro to ``quantum''.}
How to use the universal enveloping
algebra $U\mathfrak{g}$ of
$\mathfrak{g}$? Consider
$\mathfrak{g}^{*}$ as an algebraic
variety. Then
$$
\mathcal{O}(\mathfrak{g}^{*})=
\operatorname{Sym}^{\bullet}(\mathfrak{g})
$$
By the Poincar\'e-Birkhoff-Witt theorem,
$$
U\mathfrak{g}=
\operatorname{Sym}^{\bullet}\mathfrak{g}
$$
as a vector space.

\begin{remark}
\label{remark-l1-question}
{\bf Question.}
What is the relation between
$\mathcal{O}(\mathfrak{g}^{*})$ and
$U\mathfrak{g}$?
\end{remark}

Notice that $\mathfrak{g}^{*}$ is also a
Poisson algebra with bracket
$\{x,y\} =[x,y]$ for
$x,y\in \mathfrak{g}$, called the
Kirillov-Kostant-Souriau (KKS) Poisson
structure.

Now we can consider
$$
U_{\hbar}\mathfrak{g}=
\mathbb{C}[\hbar]\otimes
T\mathfrak{g}/(xy-yx-\hbar [x,y])
$$
so for $x,y\in \mathfrak{g}$ we will kind
of have
$\widetilde{x},\widetilde{y}\in
U_{\hbar}\mathfrak{g}$ such that
$$
\frac{\widetilde{x}\widetilde{y}-
\widetilde{y}\widetilde{x}}{\hbar}
\Big|_{\hbar =0}=[x,y]=\{x,y\}
$$
There is one more topological structure
that defines a group structure on
$\mathfrak{g}^{*}$:
$$
\begin{aligned}
+: \mathfrak{g}^{*}\times
\mathfrak{g}^{*} &\longrightarrow
\mathfrak{g}^{*}\\
(\varphi,\psi) &\longmapsto \varphi +\psi
\end{aligned}
$$
and what is this in terms of functions?
$$
\begin{aligned}
\mathcal{O}(\mathfrak{g}^{*} )
&\longrightarrow
\mathcal{O}(\mathfrak{g}^{*})\otimes
\mathcal{O}(\mathfrak{g}^{*}) \\
x &\longmapsto x\otimes 1+1\otimes x
\end{aligned}
$$

\medskip\noindent
{\bf Editorial note.}
In the original, the middle arrow of
axiom A5 was labelled with a coproduct
between the two copies of
$H\otimes H\otimes H\otimes H$; since
this cannot typecheck I read it as the
permutation of the two middle factors
and wrote $\operatorname{id}\otimes P
\otimes\operatorname{id}$. The map written
$p$ in the original diagrams is written
$P$ here, as in the remark on
permutations. The unlabelled arrow of A3
was labelled $\Delta$. Other repairs:
the antipode in a group maps $g$ to
$g^{-1}$ (the original had $g\mapsto g$);
$\mathcal{O}(\operatorname{GL})$ was read
as $\mathcal{O}(\operatorname{GL}_{n})$;
the relation $g^{n}=0$ in the third
example was read as $g^{n}=1$, as in
the summary above; the dual in the second
example was written $K=\mathcal{O}(G)$ and
was read as $H=\mathcal{O}(G)$; in the
action on tensor products, $u\otimes w$
was read as $v\otimes w$; the
trivial and fundamental
representations, labelled $V(?)$ and
$V(G)$, were labelled $V(0)$ and $V(1)$ to
agree with the theorem; in the
definition of the Verma module the
induction was written from
$U\mathfrak{g}$ and was read as from
$U\mathfrak{b}$, and the action of $e$ and
$f$ had no target vector, which I took as
$v_{n-1}^{\lambda}$ and
$v_{n+1}^{\lambda}$ (for
$v_{n}^{\lambda}=\frac{f^{n}}{n!}
v_{0}^{\lambda}$); ``$ev_{0}^{\lambda}$''
was completed by ``$=0$'';
``completely irreducible'' was read as
``completely reducible''. The
relation in the first example,
$\mathcal{O}(\operatorname{GL}_{n})=
\mathbb{C}\langle t_{ij}|t_{ij}=t_{ji}
\rangle$, is left as written, although it
presumably means that the $t_{ij}$
commute. The
question marks in the coproduct and counit
of $U\mathfrak{g}$, in the third example,
and in the second part of the theorem on
$\mathfrak{sl}_{2}$ are in the original.

\section{Second lecture: Poisson-Lie
groups and quantization}
\label{section-second-lecture}

\begin{definition}
\label{definition-l2-poisson-structure}
Let $A$ be a commutative algebra. A
{\it Poisson structure} on $A$ is
$\{\,,\} :A\otimes A\to A$ satisfying
\begin{enumerate}
\item $\{f,g\} =-\{g,f\}$.
\item ({\bf Leibniz identity.})
$\{fg,h\} =f\{g,h\}+\{f,h\} g$
\item ({\bf Jacobi identity.})
$\{f,\{g,h\} \} +\{g,\{h,f\} \}
+\{h,\{f,g\} \}=0$
\end{enumerate}
\end{definition}

\begin{remark}
\label{remark-l2-bivector}
If $X$ is a smooth variety, then
$\{\,,\}$ is equivalent to a
{\it bicovector}
$\Pi\in \Gamma(X,\Lambda^{2} TX)$.
\end{remark}

\begin{example}
\label{example-l2-poisson}
\begin{enumerate}
\item $\mathcal{O}(X)=\mathbb{C}[x,y]$,
with $\{x,y\} =1$ and
$\Pi=\partial_{x}\wedge\partial_{y}$.
Here $X\cong T^{*} \mathbb{A}$.
\item $X=\mathfrak{g}^{*}$,
$\mathcal{O}(X)=
\operatorname{Sym}^{\bullet}(\mathfrak{g})$,
$\{x,y\}=[x,y]$
\end{enumerate}
\end{example}

\begin{definition}
\label{definition-l2-poisson-lie-group}
A {\it Poisson-Lie group} $G$ is
\begin{enumerate}
\item A group $G$.
\item A Poisson bracket on $G$, so
$\Pi\in \Lambda^{2}TG$ such that
$$
G\times G\to G
$$
is Poisson.
\end{enumerate}
\end{definition}

\begin{exercise}
\label{exercise-l2-poisson-lie}
\begin{enumerate}
\item Show that $\Pi_{e}=0$.
\item Show that $\Pi$ induces a Lie
bracket on
$$
\xymatrix{
T^{*}_{e}G\otimes T^{*}_{e}G
\ar[r] \ar@{=}[d]
& T^{*}_{e}G \ar@{=}[d]\\
\mathfrak{g}^{*} \otimes \mathfrak{g}^{*}
\ar[r] & \mathfrak{g}^{*}
}
$$
{\bf Hint:}
$T^{*}_{e}G\cong
\mathcal{M}_{e}/\mathcal{M}^{2}_{e}$
\end{enumerate}
\end{exercise}

\begin{definition}
\label{definition-l2-lie-bialgebra}
A {\it Lie bialgebra} $\mathfrak{g}$ is a
Lie algebra with
$$
\delta:\mathfrak{g}\to
\mathfrak{g}\wedge \mathfrak{g}
$$
such that
\begin{enumerate}
\item $\delta^{*}$ is a Lie bracket on
$\mathfrak{g}^{*}$.
\item ({\bf Cocycle condition.})
$$
\delta([a,b])=[\delta(a),1\otimes b-
b\otimes 1]+[a\otimes 1+1\otimes a,
\delta(b)]
$$
\end{enumerate}
\end{definition}

\begin{proposition}
\label{proposition-l2-lie-bialgebra}
A Poisson-Lie group gives a Lie bialgebra
structure on $\mathfrak{g}=T_{e}G$.
\end{proposition}

\begin{example}
\label{example-l2-lie-bialgebras}
\begin{enumerate}
\item $G=(\mathfrak{g}^{*},+)$,
$$
\delta:T_{e}^{*} G\otimes T_{e}^{*} G
\to T^{*}_{e}G
$$
or equivalently
$$
[\,,]:\mathfrak{g}\otimes \mathfrak{g}
\to \mathfrak{g}
$$
\item $G=\operatorname{SL}_{2}$,
$\mathfrak{g}=\mathfrak{sl}_{2}$,
$\delta(e)=\frac{1}{2} e\wedge h$,
$\delta(f)=\frac{1}{2} f\wedge h$,
$\delta(h)=0$.
\end{enumerate}
\end{example}

\begin{theorem}[Drinfeld]
\label{theorem-l2-drinfeld}
If $G$ is simply connected, then there is
a bijection between the Poisson-Lie
structures on $G$ and the Lie bialgebra
structures on $\mathfrak{g}$.
(In general, not algebraic.)
\end{theorem}

Let $r\in \mathfrak{g}\otimes \mathfrak{g}$.

\begin{definition}
\label{definition-l2-cyb}
\begin{enumerate}
\item
$$
\begin{aligned}
\operatorname{CYB}:
\mathfrak{g}\otimes \mathfrak{g}&\to
\mathfrak{g}\otimes \mathfrak{g}\otimes
\mathfrak{g}\\
r&\longmapsto [r_{12},r_{13}]+
[r_{12},r_{23}]+[r_{13},r_{23}]
\end{aligned}
$$
$$
r=\sum a_{i}\otimes b_{i}\in
\mathfrak{g}\otimes \mathfrak{g}
$$
then
$$
r_{12}=\sum_{i}a_{i}\otimes b_{i}
\otimes 1,\qquad
r_{13}=\sum_{i}a_{i}\otimes 1\otimes
b_{i},\qquad
r_{23}=\sum_{i}1\otimes a_{i}\otimes
b_{i}
$$
\item The {\it classical Yang-Baxter
equation} is
$$
\operatorname{CYB}(r)=0
$$
\item A {\it classical $r$-matrix} is a
solution to the CYB equation.
\end{enumerate}
\end{definition}

Given $r\in \Lambda^{2} \mathfrak{g}$ such
that $\operatorname{CYB}(r)$ is
$\mathfrak{g}$-invariant, we can define
$$
\begin{aligned}
\delta: \mathfrak{g} &\longrightarrow
\mathfrak{g}\otimes \mathfrak{g} \\
x &\longmapsto
[r,x\otimes 1+1\otimes x]
\end{aligned}
$$

\begin{exercise}
\label{exercise-l2-bialgebra}
Check that $(\mathfrak{g},\delta)$ is a
bialgebra. Action of $\mathfrak{g}$ on
$\mathfrak{g}$ by
$$
x\cdot y=[x,y]\rightsquigarrow
\mathfrak{g}\curvearrowright
\mathfrak{g}\otimes \mathfrak{g}\otimes
\mathfrak{g}\ni\Xi
$$
and $\Xi\in\ ?\iff$. More explicitly,
[Several lines on the side board
concerning $\Xi$.]
\end{exercise}

\begin{remark}
\label{remark-l2-simple}
\begin{enumerate}
\item If $G$ is simple (e.g.
$G=\operatorname{SL}_{n},
\operatorname{SO}_{n},\ldots$), then any
$\delta$ comes from some
$r\in \mathfrak{g}\wedge \mathfrak{g}$.
And the implication $\impliedby$ is the
Whitehead lemma.

Given such $\Gamma=r\,?$, bivector,
$$
\Pi=r^{\ell}-r^{r}
$$
where the element on the left is
left-invariant and the one on the right
is right-invariant.
\end{enumerate}
\end{remark}

\begin{example}
\label{example-l2-gln}
$G=\operatorname{GL}_{n}\subset
\operatorname{Mat}_{n\times n}
(\mathbb{C})$,
$\mathcal{O}(\operatorname{GL}_{n})=
\mathbb{C}[t_{ij},\det{}^{-1}]$, and
$T\operatorname{GL}_{n}=
\operatorname{GL}_{n}\times
\operatorname{Mat}_{n\times n}
(\mathbb{C})$.

Then for $X\in \operatorname{Mat}_{n\times
n}(\mathbb{C})$, the left-invariant vector
field
$X^{\ell}:\operatorname{GL}_{n}\to
\operatorname{Mat}_{n\times n}$ with
$X^{\ell} (g)=g\cdot X$ and the
right-invariant $X^{r} (g)=X\cdot g$.

Recall:
$T=(t_{ij})\in \operatorname{End}
(\mathbb{C}^{n} \otimes
\mathcal{O}(\operatorname{GL}_{n}))$.
Given such
$r\in \mathfrak{gl}_{n}\otimes
\mathfrak{gl}_{n}$, the Poisson structure
is
$$
\{T\otimes T\}=[r,T_{1}T_{2}]
$$
So what are these $T_{1}$ and $T_{2}$?
Let us do an example for $n=2$.
$$
T_{1}T_{2}=\begin{pmatrix}
t_{11}t_{11}&t_{11}t_{12}&
t_{12}t_{11}&t_{12}t_{12}\\
t_{11}t_{21}&t_{11}t_{22}&\cdots \\
\vdots
\end{pmatrix}
$$
and so on. And then $r$ is a constant
matrix
$$
r=\begin{pmatrix}
r_{11} &\cdots &r_{1n}\\&&\\&
\end{pmatrix}
$$
\end{example}

Let us do another example.
\begin{example}
\label{example-l2-gln-r-matrix}
$G=\operatorname{GL}_{n}$,
$r=E_{12}\otimes E_{21}+
E_{11}\otimes E_{11}+
E_{22}\otimes E_{22}$. Take
$$
r=\begin{pmatrix}
1&0&0&0\\0&0&0&0\\0&0&1&0\\0&0&0&1
\end{pmatrix}
$$
So
$$
\begin{aligned}
&\begin{pmatrix}
\{t_{11},t_{11}\}&\{t_{11},t_{12}\}
&\cdots \\
\vdots &&
\end{pmatrix}\\
&=\begin{pmatrix}
1&0&0&0\\0&0&0&0\\0&0&1&0\\0&0&0&1
\end{pmatrix}
\cdot
\begin{pmatrix}
\{t_{11},t_{11}\}&\{t_{11},t_{12}\}
&\cdots \\
\vdots &&
\end{pmatrix}\\
&\quad -\begin{pmatrix}
\{t_{11},t_{11}\}&\{t_{11},t_{12}\}
&\cdots \\
\vdots &&
\end{pmatrix}
\cdot
\begin{pmatrix}
1&0&0&0\\0&0&0&0\\0&0&1&0\\0&0&0&1
\end{pmatrix}
\cdot \ldots
\end{aligned}
$$

Also consider the case of
$G=\operatorname{SL}_{n}$ with
$r=e\otimes f+\frac{1}{4}h\otimes h$
\end{example}

\begin{exercise}
\label{exercise-l2-difference}
Interpret the difference.
\end{exercise}

\begin{definition}
\label{definition-l2-quasi-triangular}
A {\it quasi-triangular Lie bialgebra} is
a triple $(\mathfrak{g},[\,,],r)$ where
\begin{itemize}
\item $(\mathfrak{g},[\,,])$ is a Lie
algebra
\item $r\in \mathfrak{g}\otimes
\mathfrak{g}$ is such that
\begin{enumerate}
\item $r+r^{21}$ is
$\mathfrak{g}$-invertible.
\item $\operatorname{CYB}(r)=0$.
\end{enumerate}
$r$ is an $r$-matrix.
\end{itemize}
\end{definition}

\begin{remark}
\label{remark-l2-symmetric-part}
In ??
$$
r^{2}r^{21}=E_{12}\otimes E_{21}+
E_{11}\otimes E_{11}+
E_{22}\otimes E_{22}+
E_{21}\otimes E_{12}
$$
is a bilinear form $(X,Y)\mapsto T\,?$
\end{remark}

\medskip\noindent
{\bf Quantization.}

\begin{definition}[``Definition'']
\label{definition-l2-quantization}
Let $(A,\{\,,\} )$ be a Poisson algebra.
A {\it quantization} of $A$ is $A_{\hbar}$
such that
\begin{enumerate}
\item $A_{\hbar}\cong
A\otimes_{\mathbb{C}}\mathbb{C}[\hbar]$ or
$A\otimes_{\mathbb{C}}
\mathbb{C}[[\hbar]]$ as vector spaces.
\item $A_{\hbar}/(\hbar)\cong A$ as
algebras.
\item For $f,g\in A$ consider
$\widetilde{f},\widetilde{g}\in
A_{\hbar}$ such that
$\widetilde{f}|_{\hbar=0}=f$ and
$\widetilde{g}|_{\hbar=0}=g$.
Consider their commutator: because of
condition 2., this commutator is
divisible by $\hbar$, so it makes sense
to ask that the Poisson bracket is
$$
\frac{\widetilde{f}\widetilde{g}-
\widetilde{g}\widetilde{f}}{\hbar}
\Big|_{\hbar =0}=\{f,g\}
$$
\end{enumerate}
\end{definition}

\begin{example}
\label{example-l2-quantizations}
\begin{enumerate}
\item $A=\mathcal{O}(T^{*}
\mathbb{A}^{1})=\mathbb{C}[x,y]$.
$$
A_{\hbar}=\mathbb{C}[\hbar]
\left\langle{}\widetilde{x},
\widetilde{y}\right\rangle{} /
(\widetilde{x}\widetilde{y}-
\widetilde{y}\widetilde{x}=\hbar)
$$
so
$$
A_{\hbar}\cong \mathbb{C}[\hbar]
\left\langle{}x,\hbar \partial_{x}
\right\rangle{}
$$
\item $U_{\hbar}\mathfrak{g}$ as a
quantization of
$(\mathfrak{g}^{*},\{\,,\} )$ where
$\{\,,\}$ is KKS.
\item $\mathcal{O}(\operatorname{GL}_{n})
=\mathbb{C}\left\langle{}t_{ij}
\right\rangle{} /(T_{1}T_{2}=T_{2}T_{1})$.
Let $r$ be the $r$-matrix (for $n=2$ as
before) and
$\{T\otimes ,T\} =[r,T_{1}T_{2}]$
\end{enumerate}
\end{example}

\begin{definition}
\label{definition-l2-frt}
Let $R$ be any element in
$\operatorname{End}(\mathbb{C}^{n})\otimes
\operatorname{End}(\mathbb{C}^{n})$. Then
we can define
\begin{equation}
\label{equation-l2-frt}
\frac{\mathbb{C}\left\langle{}t_{ij}
\right\rangle{} }{(RT_{1}T_{2}=
T_{2}T_{1}R)}
\end{equation}
which was defined by FRT.
\end{definition}

\begin{exercise}
\label{exercise-l2-frt}
The FRT algebra is a bialgebra with
$\Delta T=T\otimes T$. Where is the
quantization part? Assume that we choose
this $R$ in a very specific way:
$$
R=1+\hbar r+\mathcal{O}(\hbar^{*})
$$
where $1$ is the identity matrix. Show
that the $\hbar^{1}$ term of
Equation~\ref{equation-l2-frt} implies\ ?
\end{exercise}

\begin{definition}
\label{definition-l2-qybe}
Let $R\in \operatorname{End}
(\mathbb{C}^{n})\otimes
\operatorname{End}(\mathbb{C}^{n})$. The
{\it quantum YBE} is
$$
R_{12}R_{13}R_{23}=R_{23}R_{13}R_{12}
$$
in $\operatorname{End}
(\mathbb{C}^{n})^{\otimes 3}$.
\end{definition}

\begin{exercise}
\label{exercise-l2-qybe}
Let $R=1+\hbar r+\mathcal{O}(\hbar^{2})$.
Then $\hbar^{2}$ of QYBE implies CYBE for
$r$. Also, solutions are (quantum)
$R$-matrices.
\end{exercise}

\begin{example}
\label{example-l2-r-matrix-q}
Consider
$$
R=\begin{pmatrix}
1&0&0&0\\0&1&0&0\\0&q-q^{-1}&1&0\\0&0&0&q
\end{pmatrix}
$$
we may consider $q=e^{\hbar}$:
$$
R=1+\hbar \begin{pmatrix}
1&0&0&0\\0&0&0&0\\0&2&0&\\0&0&0&1
\end{pmatrix}
$$
$$
\hbar \cdot (E_{11}\otimes E_{11}+
E_{22}\otimes E_{22}+
2E_{12}\otimes E_{21})
$$
\end{example}

\begin{definition}
\label{definition-l2-quantum-gl2}
$\mathcal{O}(\operatorname{GL}_{2})=
\mathbb{C}\left\langle{}t_{ij}
\right\rangle{} /
(RT_{1}T_{2}=T_{2}T_{1}R)$ subject to
$\operatorname{qdet}$ is invertible,
where
$$
\operatorname{qdet}=\sum_{\sigma\in S_{n}}
(-q)^{e(\sigma)}
t_{n\,?}\ldots t_{?}
$$
\end{definition}

\medskip\noindent
{\bf Editorial note.}
Repairs: the third term of the classical
Yang-Baxter map was written
$[r_{13},r_{13}]$ and was read as
$[r_{13},r_{23}]$; the definition of
$r_{13}$ ended in the garbled
``$r_{13=\ldots}$'', which I read as the
definition of $r_{23}$; $\partial_{x\wedge
\partial_{y}}$ was read as
$\partial_{x}\wedge\partial_{y}$;
$\operatorname{End}(\mathbb{C}^{*})$ in
the definition of the FRT algebra was
read as $\operatorname{End}
(\mathbb{C}^{n})$, and the sentence
``was defined by FRT'' was attached to
the displayed algebra; $R_{1}$
in the relation of the quantum
$\mathcal{O}(\operatorname{GL}_{2})$ was
read as $R$; in the quantization the
restriction $|_{\hbar}$ was read as
$|_{\hbar=0}$; in the last example, the
terms of the displayed element, written
$E_{^{-1}}$ and $E_{2}$, were read as
$E_{11}$ and $E_{22}$. In the
Poisson-Lie theorem the original
displayed a bijection between two sets
written with text inside the formula,
which I wrote in words. The side-board
lines on $\Xi$, the remark ``In ??'', and
the formula for the quantum determinant
are incomplete in the original; the
statement of Example 3 of the
quantizations ends in the original with
a magenta comma between $T\otimes$ and
$T$, written here as a comma.

\section{Lecture 3: The quantum group of
$\mathfrak{sl}_{2}$ and link invariants}
\label{section-lecture-3}

\begin{definition}[Main definition]
\label{definition-l3-uq-sl2}
$U_{q}(\mathfrak{sl}_{2})$ is a Hopf
algebra over $\mathbb{C}(q)$ generated by
$E,F,K^{\pm 1}$ such that
\begin{enumerate}
\item $KK^{-1} =K^{-1} K=1$.
\item $KEK^{-1}=q^{2}E$.
\item $KFK^{-1} =q^{-2} F$.
\item $EF-FE=\frac{K-K^{-1}}{q-q^{-1}}$
\end{enumerate}
with coproduct
$$
\begin{aligned}
\Delta(E)&=E\otimes 1+K\otimes E\\
\Delta(F)&=1\otimes F+F\otimes K^{-1} \\
\Delta(K)&=K\otimes K
\end{aligned}
$$
counit
$$
\eta(E)=\eta(F)=0,\quad \eta(K)=1
$$
and antipode given by $S(K)=K^{-1}$,
$S(E)=-K^{-1} E$ and $S(F)=-FK$.
\end{definition}

Quantization? Take $q=e^{\hbar}$.
Informally,
$$
U_{q}(\mathfrak{sl}_{2})
\overset{\hbar\to 0}{\longrightarrow}
U(\mathfrak{sl}_{2})
$$
$$
E\mapsto e,\quad F\mapsto f,\quad
K\mapsto e^{\hbar h}
$$
$$
\frac{K-K^{-1}}{q-q^{-1}}=
\frac{e^{\hbar h}-e^{-\hbar h}}
{e^{\hbar}-e^{-\hbar}}=
\frac{\hbar(2h)}{2\hbar}
\to \mathcal{O}(\hbar) =h
$$

Let us take
$$
\frac{\Delta-\Delta^{\operatorname{op}}}
{\hbar}\Big|_{\hbar=0}
$$
What is this applied to $E$?
$$
\begin{aligned}
\frac{\Delta-\Delta^{\operatorname{op}}}
{\hbar}\Big|_{\hbar=0}(E)
&=\frac{E\otimes 1+e^{\hbar h}\otimes E
-1\otimes E-E\otimes e^{\hbar h}}{\hbar}\\
&=h\otimes e-e\otimes h=h\wedge e
\end{aligned}
$$

\begin{remark}
\label{remark-l3-upshot}
{\bf Upshot.}
This produces the coalgebra structure
from last lecture, i.e. the cobracket:
$$
\frac{\Delta-\Delta^{\operatorname{op}}}
{\hbar}\Big|_{\hbar =0}=\delta
$$
So this is a quantization in the sense of
coalgebras.
\end{remark}

\medskip\noindent
{\bf Finite-dimensional representations.}

\begin{definition}[Notations]
\label{definition-l3-quantum-integers}
\begin{enumerate}
\item A {\it quantum integer}
$$
[n]:=\frac{q^{n}-q^{-n}}{q-q^{-1}}
$$
\item The {\it quantum factorial} is
$$
[n]!:=\prod_{i=1}^{n}[i]
$$
\end{enumerate}
\end{definition}

\begin{theorem}
\label{theorem-l3-representations}
\begin{enumerate}
\item Any finite dimensional
representation of $U_{q}(\mathfrak{sl}_{2})$
is {\it completely reducible}, i.e.,
$V=\bigoplus_{i}V_{i}$ with $V_{i}$
irreducible.
\item Any irreducible representation is
isomorphic to
$V(n)=\operatorname{span}
(v_{0},\ldots,v_{n})$,
$$
\begin{aligned}
Kv_{i}&=\lambda q^{-2i} v_{i},
\qquad \lambda=q^{n}\\
Fv_{i}&=[i+1]v_{i+1}\\
Ev_{i}&=\frac{\lambda q^{i-1}-
\lambda^{-1} q^{i-1}}{q-q^{-1}}v_{i-1}
\end{aligned}
$$
\end{enumerate}
\end{theorem}

\begin{exercise}
\label{exercise-l3-check}
Check the formulas $\ddot\smile$ The point
is that $v_{i}=\frac{F^{i}}{[i]!}v_{0}$.
\end{exercise}

\begin{remark}
\label{remark-l3-infinite}
This is not necessarily
finite-dimensional.
\end{remark}

\begin{definition}
\label{definition-l3-verma-module}
The {\it Verma module} of weight
$\lambda \in \mathbb{C}(q)^{*}$ is either
\begin{itemize}
\item $M_{\lambda}:=\operatorname{span}
(v_{0}^{\lambda},\ldots,v_{u}^{\lambda})$
with the action given above, or
\item $M_{\lambda}=U_{q}\mathfrak{sl}_{2}
\otimes_{U_{q}(\mathfrak{b})}
\mathbb{C}_{\lambda}$.
\end{itemize}
\end{definition}

\begin{exercise}
\label{exercise-l3-important}
This is an important exercise.
\begin{enumerate}
\item $V(1)\otimes V(1)$ is a
$U_{q}(\mathfrak{sl}_{2})$-module.
Decompose it:
$V(1)\otimes V(1)\cong
V(2)\oplus V(0)$.

So we have $V(1)\otimes V(q)\cong
S^{1}(V)\oplus \Lambda^{2}(V)$.

\item Recall that in general
$V\otimes W\not\cong W\otimes V$ for
$V,W \in \operatorname{Rep}(H)$ of a Hopf
algebra, but, show that
$\hat{R}:=R\cdot P$ with $P$ permutations,
defines an isomorphism of
$V(1)\otimes V(1)\cong V(1)\otimes V(1)$
as
$U_{q}(\mathfrak{sl}_{2})$-representations.
\end{enumerate}
\end{exercise}

\medskip\noindent
{\bf Knot invariant.}

\begin{definition}
\label{definition-l3-knot}
A {\it knot} is a smooth oriented
embedding $S^{1}\hookrightarrow S^{3}$.
A {\it link} is a smooth oriented
embedding
$S^{1}\cup \ldots\cup S^{1}
\hookrightarrow S^{3}$.

A {\it link diagram} is a projection to
$\mathbb{R}^{2}$ such that the worst
singularities are double points.
\end{definition}

\begin{example}
\label{example-l3-links}
[Drawings do not show crossings
correctly.]
\begin{enumerate}
\item Unknot. [Figure in the original
notes: a circle.]
\item Hopf link. [Figure in the original
notes: two linked circles.]
\item Trefoil (right or left-handed).
[Figure in the original notes: a trefoil
knot drawn as a parametric curve with an
arrow.]
\end{enumerate}
\end{example}

\begin{theorem}[Lickorish-Wallace]
\label{theorem-l3-lickorish-wallace}
Any closed 3-dimensional manifold can be
constructed via a {\it Dehn surgery}.
\end{theorem}

\begin{definition}
\label{definition-l3-knot-invariant}
A {\it knot invariant} is an assignment of
a number to a link that respects isotopy.
\end{definition}

\begin{theorem}[Reidemeister,
Alexander-Briggs]
\label{theorem-l3-reidemeister}
The three {\it Reidemeister moves} are:
\begin{enumerate}
\item[(R1)] [no pictures...]
\item[(R2)]
\item[(R3)]
\end{enumerate}
Two diagrams come from isotopic links if
and only if they are related by
Reidemeister moves.
\end{theorem}

Physicists introduced the idea that
morphisms between algebras can be
represented by diagrams. So the idea is to
view links as homomorphisms of
$U_{q}(\mathfrak{sl}_{2})$-representations.
These are the rules:
\begin{enumerate}
\item Empty diagram is trivial
representation: $\varnothing
\longleftrightarrow\mathbb{C}$
\item Identity morphism is a vertical
line: $\downarrow\longleftrightarrow
\operatorname{id}_{V}$.
\item A diagram $A$ with $k$ lines coming
in and $n$ coming out (that are
representations), labelled
$V_{1},\ldots,V_{k}$ and
$W_{1},\ldots,W_{m}$, corresponds to a
homomorphism $A \in \operatorname{Hom}
(V_{1}\otimes \ldots\otimes V_{k},
W_{1}\otimes \ldots\otimes W_{n})$.

This should respect structure:
\item {\it Stacking vertically} two
diagrams $A,B$ like the one in the last
item yields a third one. This corresponds
to composition, i.e.
$B\circ A \in \operatorname{Hom}
(V_{1}\otimes \ldots\otimes V_{j},
U_{1}\otimes \ldots\otimes U_{m})$.
\item {\it Stacking horizontally}
corresponds to tensor product of
homomorphisms,
$A\otimes B \in \operatorname{Hom}
(V_{1}\otimes V_{2},W_{1}\otimes W_{2})$.
\end{enumerate}

We restrict to diagrams whose in/output
is a tensor product of an irreducible
representation $V$. Then, given a link, it
is enough to specify just three morphisms
that look like $\cup$, $\cap$ and $X$
(last one with a crossing). These are
called, respectively, {\it coevaluation},
{\it evaluation} and $\hat{R}$.

\begin{exercise}
\label{exercise-l3-cup-cap}
\begin{enumerate}
\item $U(\mathfrak{sl}_{2})$,
$V=\mathbb{C}^{2}$. Corresponding to
$\cup$ with $\mathbb{C}$ below and
$\mathbb{C}$ on each end above, we get the
representation $\mathbb{C}\to V\otimes V$
with $1\mapsto u_{0}\otimes v_{1}-
u_{1}\otimes v_{0}$ using that
$V\otimes V\cong V(1)\otimes V(1)$. This
representation is determined up to a
scalar $c_{1}$.

Another diagram $\cap$ with a copy of
$\mathbb{C}$ above and two in the lower
ends yields another representation
determined up to a scalar $c_{2}$.

[An $S$-shaped diagram corresponding to
$V\to V$, showing isotopy invariance.]

The exercise is to determine $c_{1}$ and
$c_{2}$ and their relation.

\item Do the same for
$U_{q}(\mathfrak{sl}_{2})$.
\end{enumerate}
\end{exercise}

If this is well-defined, then the third
Reidemeister move gives us
$$
(R 3)\longleftrightarrow
\check R_{12}\check R_{23}\check R_{12}=
\check R_{23}\check R_{12}\check R_{23}
$$
This is called the {\it braid relation}.
If $\check R=R\circ P$, then that is also
equivalent to
$R_{12}R_{13}R_{23}=R_{23}R_{13}R_{12}$.

\begin{theorem}
\label{theorem-l3-link-invariant}
Correspondence between links and
homomorphisms of
$U_{q}(\mathfrak{sl}_{2})$-representations
gives a link invariant.
\end{theorem}

\medskip\noindent
{\bf Editorial note.}
Repairs: in the relations of
$U_{q}(\mathfrak{sl}_{2})$ the original had
$KEK^{1}=q^{2}E$ and
$KEK^{-1}=q^{-2}F$, read as
$KEK^{-1}=q^{2}E$ and
$KFK^{-1}=q^{-2}F$; in the coproduct the
original had $\oplus$ where I wrote
$\otimes$; in the first theorem the
statement ``completely reducible, with
$V_{i}$ reducible'' was read with $V_{i}$
irreducible, and $Kv_{i}=\lambda q^{-2i}
v_{1}$ was read with $v_{i}$ on the right;
the ``Vernon model'' is the Verma module;
the third Reidemeister move, in the
original $\check R_{12}\check R_{23}
\check R_{12}=\check R_{23}\check R_{12}
\check R_{13}$, was read as the braid
relation with $\check R_{23}$ at the end;
the quotation marks in the informal limit
of $U_{q}(\mathfrak{sl}_{2})$ were turned
into the word ``informally''. Left as in
the original, although they look wrong:
the $q$-exponent in the action of $E$ on
$v_{i}$ in the theorem, the last vector
$v_{u}^{\lambda}$ of $M_{\lambda}$, the
statement $V(1)\otimes V(q)\cong
S^{1}(V)\oplus\Lambda^{2}(V)$, and the
display ending in $\mathcal{O}(\hbar)=h$.
The three pictures of an unknot, a Hopf
link and a trefoil, and the (empty)
list of Reidemeister moves, are not
reproduced.

\section{Lecture 4: Yangians}
\label{section-lecture-4}

\begin{proposition}
\label{proposition-l4-given}
Given a ...
\end{proposition}

\begin{definition}
\label{definition-l4-cybe-spectral}
Let $r(z)\in \mathfrak{g}\otimes
\mathfrak{g}((z))$. The {\it CYBE with
spectral parameter} is
$$
\begin{aligned}
\operatorname{CYB}(r(z))&=
[r_{12}(z_{1}-z_{2}),r_{13}(z_{2}-z_{3})]\\
&\quad +[r_{12}(z_{1}-z_{2}),
r_{23}(z_{2}-z_{3})]\\
&\quad +[r_{13}(z_{1}-z_{3}),
r_{23}(z_{1}-z_{3})]
\end{aligned}
$$
for $|z_{1}|<|z_{2}| <|z_{3}|$.
\end{definition}

\begin{exercise}
\label{exercise-l4-check}
Check
$\operatorname{CYB}\left(
\frac{\sum_{i}z_{i}\otimes x_{i}}{z}
\right) =0$.
\end{exercise}

\begin{theorem}[Belavin-Drinfeld]
\label{theorem-l4-belavin-drinfeld}
Let $r(z)$ satisfy
$\operatorname{CYB}(r(z))=0$ and a
non-degeneracy condition. Then
\begin{enumerate}
\item $r(z)$ extends to a meromorphic
function on $\mathbb{C}$.
\item Poles of $r(z)$ are on a lattice
$\Gamma\subset \mathbb{C}$. (E.g.
$\Gamma=\{0\},\mathbb{Z},\mathbb{Z}^{2},
\ldots$.)
\item $r(z)=\frac{\varepsilon
\sum_{i}z_{i}\otimes x_{i}}{z}
+\mathcal{O}(1)$.
\item $r(z)+r^{2?}(-z)=0$.
\end{enumerate}
\end{theorem}

$$
R(z)=1-\frac{P}{z}\in
\operatorname{End}(\mathbb{C}^{n}\otimes
\mathbb{C}^{n})[z^{-1}]
$$
$$
P=\sum_{i,j}E_{ij}\otimes E_{ji}
$$

\begin{definition}
\label{definition-l4-yangian}
The {\it Yangian}
$\operatorname{Y}(\mathfrak{gl}_{n})$ is an
algebra generated by $t_{ij}^{(k)}$,
$i,j=1,\ldots,n$ subject to
$$
t_{ij}(u):=\delta_{ij}+\sum_{k\geq 1}
t^{(k)}_{ij}u^{-k}
$$
where $\delta_{ij}\in
\operatorname{Y}(\mathfrak{gl}_{n})[[\,?\,]]$
\end{definition}

\begin{proposition}
\label{proposition-l4-qybe}
$R(u)$ satisfies the QYBE with spectral
parameter
$$
R_{12}(u_{1}-u_{2})R_{13}(u_{1}-u_{3})
$$
\end{proposition}

[Half a board missing...]

\medskip\noindent
{\bf Fact (vv).}
There is a homomorphism
$$
\begin{aligned}
\operatorname{ev}:
\operatorname{Y}(\mathfrak{gl}_{n})
&\longrightarrow U(\mathfrak{gl}_{n}) \\
t_{ij}^{(k)} &\longmapsto
\delta_{k,1}E_{ij}\\
t_{ij}(u)&\longmapsto
\delta_{ij}+\frac{E_{ij}}{u}
\end{aligned}
$$

\medskip\noindent
{\bf How to use?}
$$
\operatorname{qdet}(T(u)):=
\sum_{c \in S_{n}}(-1)^{?}
t_{ij}(u)\cdot t_{?,2}(u-1)\cdot \ldots
\cdot t_{?,n}(u-n+1)\in
\operatorname{Y}(\mathfrak{gl}_{n})
[[u^{-1}]]
$$

\begin{proposition}[v]
\label{proposition-l4-central}
$\operatorname{qdet}(T(u))$ is central.
{\bf Idea}: $R(u)=$ antisymmetrizer.
\end{proposition}

Combining the proposition (v) and the
fact (vv) we get the following.

\begin{theorem}
\label{theorem-l4-center}
$$
u\cdot (u-1)\ldots(u-n-1)\,
\operatorname{ev}(\operatorname{qdet}(T(u)))
$$
generates $Z(U(\mathfrak{gl}_{n}))$.
\end{theorem}

\medskip\noindent
{\bf Editorial note.}
Lecture 4 is unfinished in the original:
the proposition at the start, the second
proposition (the quantum Yang-Baxter
equation with spectral parameter), the
definition of the Yangian and the formula
for the quantum determinant are cut off or
carry question marks, and half a board is
noted as missing. Repairs: the equations
were called CBYE, CRB and CBY in the
original and are all called CYB here;
the formula for $R(z)$ had a stray control
sequence before the $1$ and a capital $Z$,
and was read as $1-\frac{P}{z}$; in the
homomorphism
$\operatorname{ev}$ the image of
$t_{ij}(u)$ was written
$\delta_{ij}+\frac{E_{ij}}{4}$ and was read
as $\delta_{ij}+\frac{E_{ij}}{u}$;
``$R(U)=$ antisymmetrizer'' was read with
$u$. The numerator
$\sum_{i}z_{i}\otimes x_{i}$ in the last
two items of the Belavin-Drinfeld theorem,
the factor $u(u-1)\ldots(u-n-1)$ and the
first factor $t_{ij}(u)$ in the quantum
determinant are left as in the original,
although they look wrong. The statement of
the last theorem has no condition on
$u$.


\bibliography{my}
\bibliographystyle{amsalpha}

\end{document}
